\section{Circuit Complexity --- Uniform Tableau Emitter}
%%% generated by the blueprint pipeline from TCSlib.Complexity.CircuitComplexity.UniformTableauEmitter
\subsection*{Circuit Complexity --- Uniform Tableau Emitter: Introduction}

Declarations of the Lean module \texttt{TCSlib.Complexity.CircuitComplexity.UniformTableauEmitter}.

\subsection*{Circuit Complexity --- Uniform Tableau Emitter: Results}

\begin{definition}[Execution record fields]
\lean{Complexity.UTab.ER}
\label{Complexity.UTab.ER}
\leanok
An execution record is a tuple of eleven natural numbers: the number of circuit inputs,
the number of hard-wired bits, the step budget, the current instruction, the cell used
by that instruction one step earlier, the first instruction of the current step, the
step-loop counter, the inner-loop counter, the position in the hard-wired input, the
position in the free input, and a scratch value.
\end{definition}

\begin{definition}[Execution record coordinates]
\lean{Complexity.UTab.ER.f}
\label{Complexity.UTab.ER.f}
\leanok
\uses{Complexity.UTab.ER}
For an execution record $e$, define its eleven coordinates, in order, to be its number
of circuit inputs, number of hard-wired input bits, step budget, current instruction,
preceding cell of that instruction, first instruction of the current step, step-loop
counter, inner-loop counter, position in the hard-wired input, position in the free
input, and scratch value.
\end{definition}

\begin{definition}[Updating an execution record coordinate]
\lean{Complexity.UTab.ER.set}
\label{Complexity.UTab.ER.set}
\leanok
\uses{Complexity.UTab.ER}
Given an execution record $e$, a coordinate index $r\in\{0,1,\ldots,10\}$, and a natural
number $v$, define the updated record by replacing coordinate $r$ with $v$ and leaving
every other coordinate unchanged. The coordinates are ordered as the number of circuit
inputs, number of hard-wired input bits, step budget, current instruction, preceding
cell of that instruction, first instruction of the current step, step-loop counter,
inner-loop counter, position in the hard-wired input, position in the free input, and
scratch value.
\end{definition}

\begin{theorem}[Input-count register lookup]
\lean{Complexity.UTab.ER.f_nc}
\label{Complexity.UTab.ER.f_nc}
\leanok
\uses{Complexity.UTab.ER, Complexity.UTab.ER.f, Complexity.UTab.rNC}
\difficulty{1}
For every register record $e$, the value of its register function at the input-count
index is $e.nc$.
\end{theorem}

\begin{theorem}[Hard-wired length lookup]
\lean{Complexity.UTab.ER.f_z}
\label{Complexity.UTab.ER.f_z}
\leanok
\uses{Complexity.UTab.ER, Complexity.UTab.ER.f, Complexity.UTab.rZ}
\difficulty{1}
For any register state $e$, reading its hard-wired-input-length coordinate returns the
length stored in $e$.
\end{theorem}

\begin{theorem}[Time register lookup]
\lean{Complexity.UTab.ER.f_t}
\label{Complexity.UTab.ER.f_t}
\leanok
\uses{Complexity.UTab.ER, Complexity.UTab.ER.f, Complexity.UTab.rT}
\difficulty{1}
For any register state $e$, reading its time-budget coordinate returns the time budget
stored in $e$.
\end{theorem}

\begin{theorem}[Current instruction field lookup]
\lean{Complexity.UTab.ER.f_i}
\label{Complexity.UTab.ER.f_i}
\leanok
\uses{Complexity.UTab.ER, Complexity.UTab.ER.f, Complexity.UTab.rI}
\difficulty{1}
For a register record $e$, the function that assigns its eleven fields to indices $0$
through $10$ in order takes the value $e.i$ at index $3$.
\end{theorem}

\begin{theorem}[Register Field Lookup]
\lean{Complexity.UTab.ER.f_q}
\label{Complexity.UTab.ER.f_q}
\leanok
\uses{Complexity.UTab.ER, Complexity.UTab.ER.f, Complexity.UTab.rQ}
\difficulty{1}
Let $e$ be a register record, and let $f_e$ be the function from the eleven indices
$0,\ldots,10$ to the natural numbers defined by its fields in order. At index $4$, this
function has value equal to the record's field for the current instruction's cell one
step earlier.
\end{theorem}

\begin{theorem}[Step-loop register lookup]
\lean{Complexity.UTab.ER.f_lb}
\label{Complexity.UTab.ER.f_lb}
\leanok
\uses{Complexity.UTab.ER, Complexity.UTab.ER.f, Complexity.UTab.rLB}
\difficulty{1}
Let $e$ be a register record, and let $f_e$ assign each of the eleven indices from $0$
to $10$ to the corresponding register value. At index $5$, this function returns the
step-loop counter $e.lb$.
\end{theorem}

\begin{theorem}[Loop-counter register value]
\lean{Complexity.UTab.ER.f_lt}
\label{Complexity.UTab.ER.f_lt}
\leanok
\uses{Complexity.UTab.ER, Complexity.UTab.ER.f, Complexity.UTab.rLT}
\difficulty{1}
For a register record $e$, the value of its register function at index $6$ is its inner
step-loop counter $e.lt$.
\end{theorem}

\begin{theorem}[Inner-loop counter lookup]
\lean{Complexity.UTab.ER.f_ll}
\label{Complexity.UTab.ER.f_ll}
\leanok
\uses{Complexity.UTab.ER, Complexity.UTab.ER.f, Complexity.UTab.rLL}
\difficulty{1}
For every register record $e$, the function assigning its eleven fields to indices $0$
through $10$ takes the value $e.ll$ at index $7$.
\end{theorem}

\begin{theorem}[Lookup of the hard-wired input position]
\lean{Complexity.UTab.ER.f_pa}
\label{Complexity.UTab.ER.f_pa}
\leanok
\uses{Complexity.UTab.ER, Complexity.UTab.ER.f, Complexity.UTab.rPA}
\difficulty{1}
For any register record $e$, the value of its field-assignment function at index $8$ is
$e.pa$, the position in the hard-wired part of the input.
\end{theorem}

\begin{theorem}[Free input position lookup]
\lean{Complexity.UTab.ER.f_pb}
\label{Complexity.UTab.ER.f_pb}
\leanok
\uses{Complexity.UTab.ER, Complexity.UTab.ER.f, Complexity.UTab.rPB}
\difficulty{1}
For a register record $e$, let $f_e$ denote the function from the indices $0,\ldots,10$ to its eleven natural-number fields, in their declared order. Its value at index
$9$ is $e.pb$, the position in the free part of the input.
\end{theorem}

\begin{theorem}[Scratch register lookup]
\lean{Complexity.UTab.ER.f_tmp}
\label{Complexity.UTab.ER.f_tmp}
\leanok
\uses{Complexity.UTab.ER, Complexity.UTab.ER.f, Complexity.UTab.rTMP}
\difficulty{1}
For any register state $e$, reading its scratch coordinate returns the scratch value
stored in $e$.
\end{theorem}

\begin{theorem}[Updating the register coordinate function]
\lean{Complexity.UTab.ER.update_f}
\label{Complexity.UTab.ER.update_f}
\leanok
\uses{Complexity.UTab.ER, Complexity.UTab.ER.f, Complexity.UTab.ER.set}
\difficulty{4}
Let $e$ be a record whose eleven natural-number fields, in order, are the circuit-input
count, hard-wired-bit count, step budget, current instruction, preceding cell, first
instruction of the current step, two loop counters, hard-wired and free input positions,
and scratch space. Let $f(e)$ be the function on slots $0$ through $10$ that lists these
fields in that order, and let setting slot $r$ to $v$ replace the corresponding field of
$e$. Updating $f(e)$ at any slot $r$ to the value $v$ gives the coordinate function of
the record obtained by setting slot $r$ to $v$.
\end{theorem}

\begin{definition}[Zero execution record]
\lean{Complexity.UTab.ER.zero}
\label{Complexity.UTab.ER.zero}
\leanok
\uses{Complexity.UTab.ER}
The zero execution record is the record in which all eleven fields are zero.
\end{definition}

\begin{theorem}[Coordinates of the zero record]
\lean{Complexity.UTab.ER.zero_f}
\label{Complexity.UTab.ER.zero_f}
\leanok
\uses{Complexity.UTab.ER.f, Complexity.UTab.ER.zero}
\difficulty{4}
Let $e$ be a record whose eleven natural-number fields, in order, are the circuit-input
count, hard-wired-bit count, step budget, current instruction, preceding cell, first
instruction of the current step, two loop counters, hard-wired and free input positions,
and scratch space. The coordinate function of the zero record assigns $0$ to every slot
from $0$ through $10$.
\end{theorem}

\begin{theorem}[One-step register increment]
\lean{Complexity.UTab.goes_incR}
\label{Complexity.UTab.goes_incR}
\leanok
\uses{Complexity.CounterProg.Goes, Complexity.CounterProg.goes_inc, Complexity.UTab.ER.f, Complexity.UTab.ER.set, Complexity.UTab.ER.update_f}
\difficulty{3}
Let $P$ be a counter program with eleven natural-number registers, let $x$ be a binary
input, and let $\rho$ be the register assignment encoded by an environment $e$. If the
instruction at label $\ell$ increments register $r$ and proceeds to label $\ell'$, then
from label $\ell$, input position $p$, and any current output string, the program reaches
label $\ell'$ in at most one step, with the register assignment obtained by increasing
register $r$ by one, the input position still $p$, and no additional output.
\end{theorem}

\begin{theorem}[One-step register decrement]
\lean{Complexity.UTab.goes_decR}
\label{Complexity.UTab.goes_decR}
\leanok
\uses{Complexity.CounterProg.Goes, Complexity.CounterProg.goes_dec, Complexity.UTab.ER.f, Complexity.UTab.ER.set, Complexity.UTab.ER.update_f}
\difficulty{3}
Let $P$ be a counter program with 11 registers, let $x$ be a binary input, and let $\rho$
be the register assignment encoded by an environment $e$. If the instruction at label
$\ell$ decrements register $r$ and jumps to label $\ell'$, then for every current output
string $o$, there is a run of at most one step from label $\ell$ with registers $\rho$ and
input position $p$ to label $\ell'$, with register assignment obtained by reducing register
$r$ by one (leaving it at zero if it was zero), input position still $p$, and output
still $o$.
\end{theorem}

\begin{theorem}[Reflexivity of the run relation]
\lean{Complexity.UTab.goes_refl}
\label{Complexity.UTab.goes_refl}
\leanok
\uses{Complexity.CounterProg.Goes, Complexity.CounterProg.run_zero}
\difficulty{2}
For any counter program with eleven registers, binary input, label, register assignment,
and input position, the run relation holds from that label and register assignment back
to the same label and assignment, with the same input position, no additional output,
and step bound zero.
\end{theorem}

\begin{definition}[Instruction phases]
\lean{Complexity.UTab.IPh}
\label{Complexity.UTab.IPh}
\leanok
There are seven instruction phases, denoted $p_0,z_A,z_B,u_A,u_B,e_1,e_2$.
\end{definition}

\begin{definition}[Layer assigned to each phase]
\lean{Complexity.UTab.IPh.lay}
\label{Complexity.UTab.IPh.lay}
\leanok
\uses{Complexity.UTab.IPh}
Assign a layer in $\{0,1,2\}$ to each instruction phase: $z_A$ and $z_B$ have layer $1$,
$u_A$ and $u_B$ have layer $2$, and $p_0,e_1,e_2$ have layer $0$.
\end{definition}

\begin{definition}[Universal table templates]
\lean{Complexity.UTab.Tpl}
\label{Complexity.UTab.Tpl}
\leanok
\uses{Complexity.UTab.IPh}
For each natural number $k$, define a template type with constructors for a prefix
marker, a zero-group marker carrying a Boolean, a dummy marker, a final marker, a cell
marker carrying an index in $\mathrm{Fin}(k)$, a Boolean, and a phase in
$\mathrm{Fin}(5)$, an input marker carrying a Boolean and an instruction phase, and a
snapshot marker carrying a Boolean.
\end{definition}

\begin{definition}[Input source lists by phase]
\lean{Complexity.UTab.inpPhS}
\label{Complexity.UTab.inpPhS}
\leanok
\uses{Complexity.UTab.IPh, Complexity.UTab.SSrc, Complexity.UTab.inpS}
For a natural number $W$, a Boolean $t_0$, and each phase
$	exttt{p0},	exttt{zA},	exttt{zB},	exttt{uA},	exttt{uB},	exttt{e1},	exttt{e2}$, assign
the source list obtained by the following rule. Start with the 11 entries
$(t_0,p_1,p_0,p_{-1},
eg p_0\land\neg p_{-1},\ell,n_0,n_1,n_2,c_1,c_2)$, interpreting Boolean entries as
constant sources, $\ell$ as the source for layer $\ell$, each enabled neighbor or
current source as its indexed source and each disabled one as the constant false source;
append $W$ sources indexed $0$ through $W-1$, each equal to the constant false source
when $t_0$ is true and otherwise the corresponding position source. The phase parameters
$(p_0,p_1,p_{-1},\ell)$ are respectively $(1,0,0,0)$, $(0,1,0,1)$, $(0,0,0,1)$,
$(0,1,0,2)$, $(0,0,0,2)$, $(0,1,1,0)$, and $(0,0,1,0)$.
\end{definition}

\begin{definition}[Template source lists]
\lean{Complexity.UTab.Tpl.srcs}
\label{Complexity.UTab.Tpl.srcs}
\leanok
\uses{Complexity.UTab.SSrc, Complexity.UTab.Tpl, Complexity.UTab.cellS, Complexity.UTab.inpPhS, Complexity.UTab.padS, Complexity.UTab.snapS}
For natural numbers $W,m$, associate a list of sources to each template as follows. A
cell template uses the cell-source list with its phase determining the three neighbor
flags, an input template uses the input-source list for its phase, and a snapshot
template uses the snapshot-source list; append constant-false sources to each such list
until its length is at least $m$ (adding exactly the difference when the list is
shorter), and assign the empty list to every other template. Here the cell-source list
consists of the time flag, the right-tape flag, three gated sources from each of the
left, middle, and right neighbor tracks, two gated current sources, and the position
sources; the snapshot-source list consists of the time flag, a position source at index
$W$ when the time flag is false (otherwise constant false), current sources at indices 1
and 2, two end-marker sources for each of the $k$ tracks, and the position sources.
\end{definition}

\begin{definition}[Configuration kind of a template]
\lean{Complexity.UTab.Tpl.kind}
\label{Complexity.UTab.Tpl.kind}
\leanok
\uses{Complexity.CfgKind, Complexity.UTab.Tpl}
For a template with a cell label $\tau$, its configuration kind is the nested option
indicating that label; for an input template, it is the outer option indicating an input
configuration with no cell label. Every other template has no configuration kind.
\end{definition}

\begin{definition}[Time flag of a template]
\lean{Complexity.UTab.Tpl.t0}
\label{Complexity.UTab.Tpl.t0}
\leanok
\uses{Complexity.UTab.Tpl}
The time flag of a cell, input, or snapshot template is its Boolean time parameter. The
time flag of each other template is true.
\end{definition}

\begin{definition}[Gate encoding bit list]
\lean{Complexity.UTab.gbits}
\label{Complexity.UTab.gbits}
\leanok
\uses{BoolCircuit.DAGGate, BoolCircuit.DAGGate.encode}
Given a finite list of gates, form a Boolean list by concatenating, for each gate in
order, the bit true followed by that gate's encoded representation.
\end{definition}

\begin{theorem}[Gate-list encoding by gate strings]
\lean{Complexity.UTab.encodeList_eq_gbits}
\label{Complexity.UTab.encodeList_eq_gbits}
\leanok
\uses{BoolCircuit.DAGGate, BoolCircuit.DAGGate.encode, BoolCircuit.encodeList, Complexity.UTab.gbits}
\difficulty{4}
For every finite list of directed acyclic Boolean gates, encoding the list by prefixing
each gate’s encoding with a one and terminating with a zero gives the concatenation of
those prefixed gate encodings, in list order, followed by a zero.
\end{theorem}

\begin{theorem}[Encoding preserves concatenation]
\lean{Complexity.UTab.gbits_append}
\label{Complexity.UTab.gbits_append}
\leanok
\uses{BoolCircuit.DAGGate, Complexity.UTab.gbits}
\difficulty{1}
For any two finite lists of directed acyclic Boolean gates, the Boolean string formed by
concatenating the prefixed encodings of the gates in their concatenation is the
concatenation of the corresponding Boolean strings for the two lists.
\end{theorem}

\begin{definition}[Gate-operation sequence]
\lean{Complexity.UTab.sgOps}
\label{Complexity.UTab.sgOps}
\leanok
\uses{BoolCircuit.GateKind, BoolCircuit.GateKind.encode, Complexity.CounterProg.LinE, Complexity.CounterProg.MOp, Complexity.CounterProg.bitsOps, Complexity.CounterProg.linOps}
For a gate kind and a list of linear expressions in 11 registers, form the operation
list beginning with an instruction to print true, followed by the bit-encoding
operations for the gate kind. For each linear expression in order, append an instruction
to print true, the operations for that expression, and an instruction to print false;
finish with another instruction to print false.
\end{definition}

\begin{theorem}[List encoding by concatenation]
\lean{Complexity.UTab.encodeList_eq_flatMap}
\label{Complexity.UTab.encodeList_eq_flatMap}
\leanok
\uses{BoolCircuit.encodeList}
\difficulty{4}
Let $\alpha$ be a type, let $f$ assign a finite Boolean string to each element of
$\alpha$, and let $l$ be a finite list of elements of $\alpha$. Encoding $l$ by
prefixing each encoded element with a one and terminating with a zero gives the
concatenation of the strings $1f(a)$ for $a$ in list order, followed by a zero.
\end{theorem}

\begin{theorem}[Unary output of gate instructions]
\lean{Complexity.UTab.flatMap_exec_sgOps}
\label{Complexity.UTab.flatMap_exec_sgOps}
\leanok
\uses{BoolCircuit.DAGGate.encode, BoolCircuit.GateKind, BoolCircuit.encodeNat, Complexity.CounterProg.LinE, Complexity.CounterProg.LinE.val, Complexity.CounterProg.MOp.exec, Complexity.CounterProg.flatMap_exec_bitsOps, Complexity.CounterProg.flatMap_exec_linOps, Complexity.CounterProg.linOps, Complexity.UTab.encodeList_eq_flatMap, Complexity.UTab.sgOps}
\difficulty{4}
Let $\rho$ assign a natural number to each of eleven registers, let $g$ be an AND, OR,
or NOT gate kind, and let $args$ be a finite list of linear expressions in those
registers. Executing the instruction sequence that prints a one, encodes the gate kind,
and then for each argument prints a one, the unary value of that expression, and a zero,
before printing a final zero, yields a one followed by the gate encoding whose arguments
are the values of the expressions in $args$ under $\rho$.
\end{theorem}

\begin{definition}[Tableau width]
\lean{Complexity.UTab.Wd}
\label{Complexity.UTab.Wd}
\leanok
\uses{BoolCircuit.tabW, Complexity.cfgArity, Complexity.cfgF, Complexity.cfgWidth}
For a Boolean finite-state Turing machine $M$, let $W_d$ be the tableau width obtained
by applying the tableau-width function to the configuration arity, configuration width,
and transition function of $M$.
\end{definition}

\begin{definition}[Base linear expression]
\lean{Complexity.UTab.baseE}
\label{Complexity.UTab.baseE}
\leanok
\uses{Complexity.CounterProg.LinE, Complexity.UTab.Wd, Complexity.UTab.rI, Complexity.UTab.rNC, Complexity.UTab.rZ}
For a Boolean finite-state Turing machine $M$, the base linear expression over 11
registers is the pair whose register list consists of $r_{NC}$, followed by $r_Z$,
followed by $W_d(M)$ copies of $r_I$, and whose constant term is $2+W_d(M)$.
\end{definition}

\begin{definition}[Instruction micro-operations]
\lean{Complexity.UTab.instrOps}
\label{Complexity.UTab.instrOps}
\leanok
\uses{BoolCircuit.tabFixed, Complexity.CounterProg.MOp, Complexity.UTab.SSrc.toLinE, Complexity.UTab.Tpl, Complexity.UTab.Tpl.kind, Complexity.UTab.Tpl.srcs, Complexity.UTab.baseE, Complexity.UTab.sgOps, Complexity.cfgArity, Complexity.cfgF, Complexity.cfgWidth, Complexity.snapWidth}
For a finite binary Turing machine $M$ and a template $	heta$, the instruction is a list
of micro-operations on eleven registers. It concatenates, first, the serialized AND
gates associated with the sources of $	heta$, and then the serialized gates in the fixed
transition table for the kind of $	heta$, with each table-gate argument represented by
the base linear expression shifted by that argument.
\end{definition}

\begin{definition}[Template programs]
\lean{Complexity.UTab.tmpl}
\label{Complexity.UTab.tmpl}
\leanok
\uses{BoolCircuit.constGate, Complexity.CounterProg.MOp, Complexity.CounterProg.bitsOps, Complexity.CounterProg.linOps, Complexity.UTab.SSrc.toLinE, Complexity.UTab.Tpl, Complexity.UTab.Wd, Complexity.UTab.gbits, Complexity.UTab.instrOps, Complexity.snapWidth}
For a finite binary Turing machine $M$, each template is assigned a list of
micro-operations on eleven registers. The preamble encodes the two constant gates false
and true; a zero-guard template encodes its specified constant gate; and the dummy
template encodes a list of $W_d$ false constant gates, where $W_d$ is the fixed table
width for $M$. The final template prints a delimiter, the linear encoding of the current
snapshot at position $	ext{snapWidth}(M)$, and another delimiter; every other template
is assigned the instruction micro-operations associated with that template.
\end{definition}

\begin{definition}[Maximum template length]
\lean{Complexity.UTab.Lmax}
\label{Complexity.UTab.Lmax}
\leanok
\uses{Complexity.UTab.Tpl, Complexity.UTab.tmpl}
For a finite binary Turing machine $M$, let $L_{\rm max}$ be the maximum length of the
micro-operation lists assigned to its templates.
\end{definition}

\begin{theorem}[Template length bound]
\lean{Complexity.UTab.length_tmpl_le}
\label{Complexity.UTab.length_tmpl_le}
\leanok
\uses{Complexity.UTab.Lmax, Complexity.UTab.Tpl, Complexity.UTab.tmpl}
\difficulty{1}
Fix a finite-tape Turing machine $M$ and a template $	heta$ for its number of work
tapes. The length of the micro-operation list generated by $	heta$ is at most the
maximum of these lengths over all templates for $M$.
\end{theorem}

\begin{definition}[Copy-state tags]
\lean{Complexity.UTab.CpS}
\label{Complexity.UTab.CpS}
\leanok
For $k\in\mathbb N$, define a finite set of tags consisting of $cA(\tau,t_0)$ and
$cB(\tau,t_0)$ for each tape index $\tau<k$ and bit $t_0$, together with $iz(t_0)$,
$iu(t_0)$, and $lb(t_0)$ for each bit $t_0$, and a single tag $lt$.
\end{definition}

\begin{definition}[Clause-state tags]
\lean{Complexity.UTab.ClS}
\label{Complexity.UTab.ClS}
\leanok
Define a finite set of tags consisting of $pa(t_0)$, $pb(t_0)$, and $lb(t_0)$, each
indexed by a bit $t_0$.
\end{definition}

\begin{definition}[Head-state tags]
\lean{Complexity.UTab.Hd}
\label{Complexity.UTab.Hd}
\leanok
For $k\in\mathbb N$, define a finite set of tags consisting of $cA(\tau,t_0)$ and
$cB(\tau,t_0)$ for each tape index $\tau<k$ and bit $t_0$, together with $iz(t_0)$ and
$iu(t_0)$ for each bit $t_0$, and a single tag $lay$.
\end{definition}

\begin{definition}[Program-label tags]
\lean{Complexity.UTab.Lb}
\label{Complexity.UTab.Lb}
\leanok
\uses{Complexity.UTab.ClS, Complexity.UTab.CpS, Complexity.UTab.Hd, Complexity.UTab.Tpl}
For $k,L\in\mathbb N$, define a finite set of program-label tags: eight unindexed tags
$rN,nInc,nOne,nOut0,rT,tInc,rZ,zInc$; tags $lay(t_0)$ indexed by a bit; tags
$tp(\theta,pc)$ with a template $\theta$ and a program counter $pc<L+1$; tags
$pI(\theta)$ and $pQ(\theta)$ for each template; tags $cp(s,ph)$ with a copy-state tag
$s$ and $ph<7$; tags $dA(\tau,t_0)$ and $dB(\tau,t_0)$ with $\tau<k$ and a bit; tags
$hd(h)$ and $dl(h)$ with a head-state tag $h$; tags
$zb(t_0),ub(t_0),ub2(t_0),iPA(t_0),iPB(t_0),eT(t_0),eT2(t_0)$ indexed by a bit; tags
$cl(s,ph)$ with a clause-state tag $s$ and $ph<2$; and a single tag $stop$.
\end{definition}

\begin{definition}[Template program counter zero]
\lean{Complexity.UTab.tpL}
\label{Complexity.UTab.tpL}
\leanok
\uses{Complexity.UTab.Lb, Complexity.UTab.Tpl}
For a template $\theta$ and natural numbers $k,L$, the label $tpL(\theta)$ is the
program-label tag $tp(\theta,0)$.
\end{definition}

\begin{definition}[Cell start template]
\lean{Complexity.UTab.cellStart}
\label{Complexity.UTab.cellStart}
\leanok
\uses{Complexity.UTab.Lb, Complexity.UTab.tpL}
Let $k,L$ be natural numbers, let $M$ be a finite Boolean Turing machine, let $t_0$ be a
Boolean, and let $\tau$ be a natural number. The cell-start label is the template for
cell $\tau$ with bit $t_0$ and phase $0$ when $\tau<k$; when $\tau\ge k$, it is the
input template with bit $t_0$ and phase $p_0$.
\end{definition}

\begin{definition}[Source register of a counter state]
\lean{Complexity.UTab.CpS.src}
\label{Complexity.UTab.CpS.src}
\leanok
\uses{Complexity.UTab.CpS, Complexity.UTab.rI, Complexity.UTab.rNC, Complexity.UTab.rT, Complexity.UTab.rZ}
Let $k$ be a natural number and let $M$ be a finite Boolean Turing machine. For a
counter state, its source register is register $2$ in states of type $cA$, $cB$, or
$lt$; register $1$ in a state of type $iz$; register $0$ in a state of type $iu$; and
register $3$ in a state of type $lb$.
\end{definition}

\begin{definition}[Destination register of a counter state]
\lean{Complexity.UTab.CpS.dst}
\label{Complexity.UTab.CpS.dst}
\leanok
\uses{Complexity.UTab.CpS, Complexity.UTab.rLB, Complexity.UTab.rLL, Complexity.UTab.rLT}
Let $k$ be a natural number and let $M$ be a finite Boolean Turing machine. The
destination register of a counter state is register $5$ for a state of type $lb$,
register $6$ for a state of type $lt$, and register $7$ for every other counter state.
\end{definition}

\begin{definition}[Counter-state exit label]
\lean{Complexity.UTab.CpS.exit}
\label{Complexity.UTab.CpS.exit}
\leanok
\uses{Complexity.UTab.CpS, Complexity.UTab.Lb}
Let $k,L$ be natural numbers and let $M$ be a finite Boolean Turing machine. A counter
state of type $cA(\tau,t_0)$ exits to the $dA$ label with parameters $\tau,t_0$, and one
of type $cB(\tau,t_0)$ exits to the corresponding $dB$ label. States $iz(t_0)$ and
$iu(t_0)$ exit to the $hd$ labels carrying those states. A state $lb(t_0)$ exits to the
counter label for $lt$ at phase $0$ if $t_0$ is true, and otherwise to the $hd$ label
for $lay$; a state of type $lt$ exits to the $hd$ label for $lay$.
\end{definition}

\begin{definition}[Register of a cleanup state]
\lean{Complexity.UTab.ClS.reg}
\label{Complexity.UTab.ClS.reg}
\leanok
\uses{Complexity.UTab.ClS, Complexity.UTab.rLB, Complexity.UTab.rPA, Complexity.UTab.rPB}
Let $k$ be a natural number and let $M$ be a finite Boolean Turing machine. The register
associated to a cleanup state is register $8$ for a state of type $pa$, register $9$ for
one of type $pb$, and register $5$ for one of type $lb$.
\end{definition}

\begin{definition}[Cleanup-state exit label]
\lean{Complexity.UTab.ClS.exit}
\label{Complexity.UTab.ClS.exit}
\leanok
\uses{Complexity.UTab.ClS, Complexity.UTab.Lb, Complexity.UTab.tpL}
Let $k,L$ be natural numbers and let $M$ be a finite Boolean Turing machine. A cleanup
state of type $pa(t_0)$ exits to the cleanup label for $pb(t_0)$ at phase $0$; one of
type $pb(t_0)$ exits to the snapshot template with bit $t_0$ at program counter $0$; and
one of type $lb(t_0)$ exits to the counter label for $lb(t_0)$ at phase $0$.
\end{definition}

\begin{definition}[Register of a head state]
\lean{Complexity.UTab.Hd.reg}
\label{Complexity.UTab.Hd.reg}
\leanok
\uses{Complexity.UTab.Hd, Complexity.UTab.rLL, Complexity.UTab.rLT}
Let $k$ be a natural number and let $M$ be a finite Boolean Turing machine. The register
associated to a head state is register $6$ for the state $lay$ and register $7$ for
every other head state.
\end{definition}

\begin{definition}[Head-state exit label]
\lean{Complexity.UTab.Hd.exit}
\label{Complexity.UTab.Hd.exit}
\leanok
\uses{Complexity.UTab.Hd, Complexity.UTab.Lb, Complexity.UTab.tpL}
Let $k,L$ be natural numbers and let $M$ be a finite Boolean Turing machine. A head
state of type $cA(\tau,t_0)$ exits to the cell template with parameters $\tau,t_0$ and
phase $2$, while one of type $cB(\tau,t_0)$ exits to the corresponding cell template at
phase $4$. A state $iz(t_0)$ exits to the counter label for $iu(t_0)$ at phase $0$, and
a state $iu(t_0)$ exits to the $eT$ label with bit $t_0$. The state $lay$ exits to the
final template at program counter $0$.
\end{definition}

\begin{definition}[Head-state body labels]
\lean{Complexity.UTab.Hd.body}
\label{Complexity.UTab.Hd.body}
\leanok
\uses{Complexity.UTab.Hd, Complexity.UTab.Lb, Complexity.UTab.tpL}
For a natural number $k$ and a finite binary Turing machine $M$, define a map from head
states to program labels as follows: a state $cA(\tau,t_0)$ maps to the template label
for cell $(\tau,t_0)$ at phase $1$, and $cB(\tau,t_0)$ maps to the corresponding
phase-$3$ label. A state $iz(t_0)$ maps to $zb(t_0)$, a state $iu(t_0)$ maps to
$ub(t_0)$, and the layout state maps to the label $lay(false)$.
\end{definition}

\begin{definition}[Template successor labels]
\lean{Complexity.UTab.Tpl.next}
\label{Complexity.UTab.Tpl.next}
\leanok
\uses{Complexity.UTab.Lb, Complexity.UTab.Tpl}
For a natural number $k$ and a finite binary Turing machine $M$, define the successor
label of a template: the prefix template maps to $rZ$, a zero-generating template maps
to $zInc$, the dummy template maps to $lay(true)$, and the final template maps to the
stop label. Every other template maps to its phase-$0$ input label.
\end{definition}

\begin{definition}[Template continuation labels]
\lean{Complexity.UTab.Tpl.after}
\label{Complexity.UTab.Tpl.after}
\leanok
\uses{Complexity.UTab.Lb, Complexity.UTab.Tpl, Complexity.UTab.cellStart}
For natural numbers $k,L$ and a finite binary Turing machine $M$, define the label
following a template as follows. A cell template at phase $0,1,2,$ or $3$ maps,
respectively, to the phase-$0$ copy label for $cA$, the head label for $cA$, the
phase-$0$ copy label for $cB$, or the head label for $cB$, retaining its tape index and
bit. A cell template at any phase at least $4$ maps to the phase-$0$ template for the
next cell, or to the phase-$0$ input template if that next cell index is outside the $k$
work tapes. An input template at phase $p0$ maps to the phase-$0$ copy label for $iz$;
at $zA$ or $zB$ it maps to $iPA$; at $uA$ or $uB$ it maps to $iPB$; and at $e1$ or $e2$
it maps to the phase-$0$ cleanup label for $pa$. A snapshot template maps to the
phase-$0$ cleanup label for $lb$, and every other template maps to the stop label.
\end{definition}

\begin{definition}[Universal simulation counter program]
\lean{Complexity.UTab.prog}
\label{Complexity.UTab.prog}
\leanok
\uses{Complexity.CounterProg.Instr, Complexity.CounterProg.tmplInstr, Complexity.UTab.ClS.exit, Complexity.UTab.ClS.reg, Complexity.UTab.CpS.dst, Complexity.UTab.CpS.exit, Complexity.UTab.CpS.src, Complexity.UTab.Hd.body, Complexity.UTab.Hd.exit, Complexity.UTab.Hd.reg, Complexity.UTab.Lb, Complexity.UTab.Lmax, Complexity.UTab.Tpl.after, Complexity.UTab.Tpl.next, Complexity.UTab.Tpl.t0, Complexity.UTab.cellStart, Complexity.UTab.rI, Complexity.UTab.rLL, Complexity.UTab.rNC, Complexity.UTab.rPA, Complexity.UTab.rPB, Complexity.UTab.rQ, Complexity.UTab.rT, Complexity.UTab.rTMP, Complexity.UTab.rZ, Complexity.UTab.tmpl, Complexity.UTab.tpL}
Let $M$ be a finite binary Turing machine with $k$ work tapes, and let $L$ be a natural
number. Define a counter program with $11$ registers and labels indexed by the templates
and control states of the simulation, using $L_{\max}(M)$ as the template
program-counter bound. Its instructions implement the following transitions: $rN$ reads
the next input bit, halting at end of input, otherwise going to $nOut0$ on $0$ and
$nInc$ on $1$; $nInc$ increments register $rNC$ and goes to $nOne$, which prints $1$ and
returns to $rN$, while $nOut0$ prints $0$ and goes to $rT$. The label $rT$ reads a bit,
halting at end of input, going to the prefix template on $0$, or to $tInc$ on $1$;
$tInc$ increments $rT$ and returns to $rT$. The label $rZ$ reads a bit, going to the
dummy template at end of input and otherwise to the zero-generating template for the bit
read; $zInc$ increments $rZ$ and returns to $rZ$. At $lay(t_0)$, the program halts if
$rT=0$ and otherwise starts processing cell $0$ with flag $t_0$. At a template label, it
executes that template’s micro-operations in order and then continues at its successor
label; the labels $pI$ and $pQ$ increment $rI$, then either continue directly when the
template flag is true or increment $rQ$ before continuing. The copy labels copy their
selected source register into the destination by decrementing the source while
incrementing the destination, using $rTMP$ to restore the source; their exit depends on
the copy state. Labels $dA$ and $dB$ decrement $rLL$ and enter the corresponding head
state. A head state branches according to whether its associated register is zero;
otherwise the program decrements that register and executes the head-state body. The
labels $zb$, $ub$, $ub2$, $iPA$, and $iPB$ test or increment the indicated registers to
select the corresponding input-template phase. The labels $eT$ and $eT2$ test $rZ$ and
$rNC$ to select an input-template phase. Cleanup labels repeatedly test and decrement
their selected register until it is zero, then continue at the cleanup state’s exit. The
stop label halts.
\end{definition}
